%% ============================================================
%%  Conclusiones preliminares
%%  Duenio inicial: claude-1
%%  Reglas: ver ../../CLAUDE.md
%%  No editar si TASKS.md marca este archivo IN_PROGRESS por otro agente.
%% ============================================================

\chapter{Conclusiones preliminares}

La revisión sistemática deja una conclusión bastante clara sobre el estado del campo: obtener una estructura relacional desde un modelo profundo está resuelto por varias vías, y decidir si esa estructura merece crédito no lo está. Hay trabajos que aprenden el grafo como parámetro, otros que lo reciben como conocimiento previo y otros que lo obtienen interviniendo las activaciones del modelo. Lo que no aparece en el corpus revisado es un procedimiento para sostener una estructura leída de un modelo ya entrenado cuando no se dispone de una referencia externa contra la cual compararla.

Sobre esa brecha, el aporte de esta etapa es un diseño completo y no un resultado. Las tres transformaciones del pipeline están formalizadas, la condición de admisibilidad que debe cumplir la matriz de entrada está explícita, la cadena de cinco eslabones está especificada con umbrales declarados antes de ejecutarla, y la taxonomía que separa relación genuina de atajo predictivo, artefacto y ruido está formulada. Corresponde ser literal en este punto: nada de eso ha sido ejecutado todavía sobre el caso de estudio, y por lo tanto este informe no reporta resultados experimentales. La ejecución es el contenido de la etapa siguiente y está calendarizada en el capítulo \ref{cap:plan}.

Del propio corpus se desprenden tres advertencias que conviene dejar anotadas antes de ejecutar el plan. La primera es que buena parte de la ganancia atribuida a los Transformers en pronóstico de largo plazo podría provenir del mecanismo de parcheo y no de la atención entre variables \cite{patchmlp2025unlocking}; la propuesta no queda invalidada, porque no reclama que la atención mejore el desempeño, pero sí queda obligada a no presentarla como fuente privilegiada de conocimiento sin mostrarlo. La segunda es que el único caso del corpus donde la atención recupera un mecanismo verificable lo hace por coincidencia con una frecuencia que la teoría ya predecía \cite{rotorcraft2026vortexring}, es decir, como hallazgo puntual y no como procedimiento repetible. La tercera es sobre la referencia externa: el trabajo más exigente en validación dispone de la topología física real de la instalación que modela \cite{pipeline2026causalgraph}, y entre índices espectrales no existe nada equivalente, con lo cual el conjunto sintético de estructura conocida deja de ser un complemento y pasa a ser la única referencia externa del plan.

Queda además una limitación del propio protocolo de revisión, que se declara aquí porque afecta la trazabilidad del corpus: el paso que va de los registros elegibles a las referencias finalmente incluidas se resolvió por juicio de relevancia documentado caso por caso, sin umbral cuantitativo explícito, y formalizarlo queda pendiente. En una línea parecida, el corpus documenta que dos configuraciones que coinciden en la predicción pueden diferir en la explicación \cite{bogaert2026variability}, lo que sugiere reportar la variabilidad entre semillas como parte del resultado y no tratarla como ruido a suprimir.
